% !TEX root = ../main.tex
\chapter{INTRODUCTION}
\label{chap:introduction}

Replace this text with your introduction. Start with the question or problem
your work addresses and explain why it matters. You can write ordinary sentences
here. A blank line starts a new paragraph.

This starter uses a book about writing with LaTeX as a real citation
example~\cite{lamport1994}. Replace it with literature relevant to your work.
Chapter~\ref{chap:research} contains small examples you can adapt, and
Chapter~\ref{chap:conclusion} closes the document.

\section{Research question}
State your central question, argument, or objective. Your project might involve
texts, archives, design, theory, experiments, interviews, or computation.
Choose chapter titles and a structure appropriate to your discipline.

\section{Context and contribution}
Explain what is already known and what your work adds. You can emphasize a
word with \emph{italics} or use \textbf{bold text} when appropriate.

If you type a percent sign, write \verb|\%|: for example, 25\%.
Write \verb|\&| for an ampersand, such as research \& development.
Other special characters are explained in the writing guide.
